\customSection[sec:analysis]{Аналіз предметної області та огляд літератури}{0em}

Перш ніж переходити до вдосконалення системи, доцільно розглянути сучасний стан предметної
області: як штучний інтелект застосовують у зміні поведінки, якими засобами збирають дані про
активність користувача та які підходи використовують наявні рішення.
Це дає змогу обґрунтовано визначити напрями розвитку системи Dailu й уникнути повторення вад
аналогів.

\customSubSection[sec:analysis_habit]{Звичка та роль зворотного зв'язку у її формуванні}{0em}

Звичка — це засвоєна модель поведінки, яка з часом виконується майже автоматично, без свідомого
вольового зусилля.
В основі більшості сучасних трекерів лежить уявлення про «петлю звички», що складається із
сигналу, дії та винагороди, а також психологічний ефект відчутності прогресу, коли видимий
ланцюжок виконаних днів мотивує користувача не переривати його.

З погляду обробки даних принциповим є те, що сама фіксація факту виконання ще не є винагородою
чи корисним висновком.
Дослідники зміни поведінки наголошують, що формування нових звичок потребує постійного
коучингу, моніторингу та підтримки, а з розвитком недорогих пристроїв моніторингу та засобів
штучного інтелекту стало звичним застосовувати ці технології для цілодобової підтримки зміни
поведінки \citecustom{4}.
Таким чином, ефективний трекер має не лише збирати дані, а й повертати користувачеві осмислений
зворотний зв'язок — саме тут відкривається простір для застосування штучного інтелекту.

\customSubSection[sec:analysis_ai]{Штучний інтелект у системах зміни поведінки}{1em}

Застосування штучного інтелекту в системах моніторингу поведінки має тривалу передісторію.
Ще на початку 2000-х років було введено поняття «переконливих технологій» (persuasive
technology) як цифрових застосунків, що підштовхують користувача до здорового способу життя;
сучасне ж покоління інструментів, представлене глибоким навчанням і великими мовними моделями,
знову підняло інтерес до цифрового здоров'я та зміни поведінки \citecustom{4}.
Огляд галузі показує, що поле поступово переходить від експериментальних розробок до реального
впровадження, охоплюючи як генеративний, так і предиктивний штучний інтелект.

У контексті трекінгу звичок засоби штучного інтелекту застосовують здебільшого у трьох формах.
По-перше, це предиктивна аналітика та рекомендації: моделі аналізують поведінкові патерни
користувача — коли й наскільки часто виконується звичка — і надають адаптивні поради в режимі,
наближеному до реального часу \citecustom{3}.
По-друге, це діалогові асистенти на основі великих мовних моделей, що виконують роль
мотиваційного коуча й формулюють персоналізовані підказки природною мовою.
По-третє, це інтелектуальна обробка й упорядкування даних — зокрема автоматична класифікація та
категоризація записів, яка знімає з користувача рутинну роботу з упорядкування власних звичок.

Аналіз наявних комерційних застосунків підтверджує цю тенденцію.
Так, окремі трекери позиціонують себе як рішення, що аналізують звички користувача та надають
персоналізовані рекомендації й «розумніші» рутини, побудовані на основі штучного інтелекту, а
також детальну аналітику прогресу \citecustom{5}.
Водночас у більшості випадків інтелектуальні функції реалізовано поверх ручного введення даних,
тобто вони не усувають головної вади — залежності обліку від дисциплінованості користувача.

\customSubSection[sec:analysis_data]{Збір даних про активність: програмні та фізичні джерела}{1em}

Окремий напрям досліджень — способи збору даних про поведінку користувача.
У межах персональної інформатики (personal informatics) — систем, що допомагають людині збирати
особисто значущі дані задля саморефлексії та самопізнання \citecustom{6}, — розрізняють два
типи збору даних.
Пасивний збір передбачає автоматичне отримання інформації без участі користувача й уможливлює
безперервний моніторинг, тоді як активний збір вимагає свідомої дії користувача в момент
фіксації \citecustom{7}.

Це розрізнення безпосередньо стосується проєкту Dailu.
Наявна інтеграція з GitHub і Strava є прикладом пасивного програмного збору: активність
фіксується зовнішніми сервісами й підтягується автоматично.
Проте значна частина звичок не залишає цифрового сліду в програмних сервісах (наприклад,
«випити склянку води», «зробити розтяжку»), і для них досі лишається тільки ручне відмічання.
Тут у нагоді стають носимі та вбудовані пристрої Інтернету речей.

Дослідження підтверджують потенціал таких пристроїв: смарт-годинники та інші носимі
IoT-пристрої містять сенсори, що підтримують збір і опрацювання даних задля висновків, які
мотивують користувача досягати цілей фізичної активності та здоров'я \citecustom{8}.
Водночас у сфері цифрового здоров'я зростає увага до простих інструментів активного збору:
описано, зокрема, «трекер з однією кнопкою» (One Button Tracker) — носимий пристрій, що дає
змогу активно, у сам момент події, фіксувати самостійно визначене, особисто значуще явище одним
натисканням \citecustom{7}.
Такий підхід демонструє, що навіть мінімалістичний фізичний пристрій може ефективно доповнювати
програмний збір даних, охоплюючи ті звички, які інакше залишилися б поза автоматизацією.

\customSubSection[sec:analysis_analogs]{Огляд наявних рішень}{1em}

Для визначення місця вдосконаленої системи розглянуто кілька категорій рішень: спеціалізовані
трекери звичок, сервіси автоматичного збору активності та застосунки з інтелектуальними
функціями.

\textbf{Habitica}, \textbf{Streaks}, \textbf{Habitify}, \textbf{Loop Habit Tracker} — популярні
спеціалізовані трекери.
Вони дають користувачеві свободу формулювати власні цілі та підтримувати щоденні ланцюжки,
проте облік у них залишається переважно ручним.
Streaks частково використовує дані Apple Health, однак ця автоматизація обмежена сферою
здоров'я й замкнена в екосистемі Apple; інтеграцій з інструментами розробника ці застосунки не
мають.

\textbf{Граф внесків GitHub} та \textbf{Strava} автоматично фіксують активність користувача
(коміти й тренування відповідно), але прив'язані лише до власної вузької предметної області та
не дають змоги визначати довільні звички чи цілі.

\textbf{Застосунки з AI-функціями} (Habi AI та подібні) аналізують поведінкові патерни й
надають персоналізовані рекомендації та мотиваційний контент \citecustom{5}.
Їхня слабка сторона — інтелектуальні функції побудовано поверх ручного введення, а джерела
даних не поєднують цифрову й фізичну активність в одній системі.

Зведене порівняння розглянутих рішень за ключовими критеріями наведено в \tabref{tab:analogs}.

\begin{center}
    \begin{xltabular}{\textwidth}{| >{\raggedright\arraybackslash}X | c | c | c | c | c |}
        \caption{Порівняння наявних рішень за ключовими критеріями} \label{tab:analogs} \\
        \hline
        \textbf{Рішення} & \textbf{Веб} & \textbf{Автозбір} & \textbf{Свої звички} & \textbf{AI-аналіз} & \textbf{IoT} \\ \hline
        \endfirsthead

        \hline
        \textbf{Рішення} & \textbf{Веб} & \textbf{Автозбір} & \textbf{Свої звички} & \textbf{AI-аналіз} & \textbf{IoT} \\ \hline
        \endhead

        Habitica & частково & ні & так & ні & ні \\ \hline
        Streaks / Habitify & ні & частково & так & частково & частково \\ \hline
        Loop Habit Tracker & ні & ні & так & ні & ні \\ \hline
        Граф внесків GitHub & так & так & ні & ні & ні \\ \hline
        Strava & так & так & ні & частково & так \\ \hline
        Habi AI (та подібні) & так & ні & так & так & ні \\ \hline
        \textbf{Dailu (розроблювана)} & так & так & так & так & так \\ \hline
    \end{xltabular}
\end{center}

\customSubSection[sec:analysis_rationale]{Обґрунтування доцільності розробки}{1em}

Зіставлення розглянутих рішень виявляє чітку прогалину.
Спеціалізовані трекери дають гнучкість у визначенні звичок, але покладаються на ручний облік і
не мають глибокого інтелектуального аналізу.
Сервіси автоматичного збору (GitHub, Strava) об'єктивно фіксують активність, проте замкнені у
власній царині.
AI-орієнтовані застосунки надають рекомендації, але не поєднують у собі автоматичний збір даних
із різнорідних джерел — і, що суттєво, жодне з розглянутих рішень не охоплює одночасно
програмну та фізичну активність з подальшим інтелектуальним опрацюванням.

Саме на заповнення цієї прогалини й спрямована робота.
Її ключова ідея — надати користувачеві гнучкість звичайного трекера, зняти з нього тягар
ручного обліку там, де активність фіксується автоматично (як програмними сервісами, так і
фізичним IoT-пристроєм), і перетворити зібрані дані на персоналізовані висновки засобами
штучного інтелекту.
Така комбінація усуває головну ваду наявних трекерів — залежність від дисциплінованості
користувача — і водночас підвищує цінність системи за рахунок інтелектуального зворотного
зв'язку.

\customSubSection[sec:analysis_requirements]{Постановка задачі та вимоги до системи}{1em}

Спираючись на результати аналізу, задачу можна сформулювати так: необхідно вдосконалити наявний
вебзастосунок Dailu, доповнивши його підсистемою інтелектуального аналізу на основі штучного
інтелекту та каналом збору фізичної активності через пристрій Інтернету речей, зберігши при
цьому наявну модульну архітектуру й можливість подальшого розширення.

\customSubSubSection[sec:req_functional]{Функціональні вимоги}{0em}

\begin{enumerate}[noitemsep, topsep=2pt, leftmargin=*, label=\arabic*)]
    \item збереження базових функцій наявної системи: реєстрація й автентифікація, керування звичками та записами, інтеграція з GitHub і Strava, візуалізація активності;
    \item генерування персоналізованих рекомендацій та інсайтів на основі історії виконання звичок користувача;
    \item автоматична категоризація звичок засобами штучного інтелекту;
    \item діалоговий асистент на основі великої мовної моделі для мотивації та порад із контекстом даних користувача;
    \item приймання подій активності від IoT-пристрою та створення відповідних записів звичок;
    \item безпечна прив'язка IoT-пристрою до облікового запису користувача.
\end{enumerate}

\customSubSubSection[sec:req_nonfunctional]{Нефункціональні вимоги}{1em}

\begin{enumerate}[noitemsep, topsep=2pt, leftmargin=*, label=\arabic*)]
    \item безпека — захищене зберігання облікових даних, токенів і ключів пристроїв, автентифікація на основі токенів;
    \item розширюваність — можливість додавати нові інтелектуальні функції та джерела активності без істотного переписування коду;
    \item підтримуваність — чіткий поділ системи на шари й модулі з вираженими межами відповідальності;
    \item надійність — стійкість фонових процесів та зовнішніх викликів (зокрема до AI-провайдера) до тимчасових збоїв;
    \item приватність — контроль користувача над обсягом даних, що передаються до AI-підсистеми;
    \item спостережуваність — структуроване журналювання й трасування роботи системи.
\end{enumerate}

Зведений перелік варіантів використання вдосконаленої системи, що охоплює як успадковані
функції, так і нові можливості на основі штучного інтелекту та IoT-пристрою, наведено на
\textbf{рис.~\ref{fig:usecase}}.

\setfigure[fig:usecase]{0.9\textwidth}{assets/diagram-usecase}{Діаграма варіантів використання вдосконаленої системи}

Сформульовані вимоги стають відправною точкою для проєктування архітектури вдосконаленої
системи, якому присвячено наступний розділ.
